\comoleerla{dos rectas, el costo total de cada ruta (peaje más km por costo variable por km) contra el costo variable por km; la ruta por Bogotá es la azul continua y el desvío por Santander la naranja discontinua; el fondo azul es la zona donde gana Bogotá y el rayado, la zona donde gana el desvío; el punto es el umbral y la línea punteada, el combustible solo.}
{las rectas se cruzan en 204.5 COP/km (un ahorro de peaje de 70\,600 COP entre 345.15 km extra); el combustible solo cuesta 1\,157.0 COP/km, a la derecha del umbral.}
{Medido: con un costo variable mayor que el umbral gana la ruta por Bogotá, y los otros costos por km (supuesto del equipo, 0 COP/km) solo pueden subirlo; no se incluyen tiempo ni riesgo.}
